% !TEX root = template-sojo-shushi-lua.tex
% !TEX recipe = LuaLaTeX (kian-lua)
% ============================================================
%  訴状（請求の趣旨）テンプレート  ― IT化対応版  (kian-v08 LuaLaTeX 版)
%  用途: 裁判IT化提出用（表紙なし・請求の趣旨のみの単独PDF）
%  ※ 表紙あり一体型は template-sojo-lua.tex を使用
%
%  コンパイル: latexmk -r latexmkrc-lua template-sojo-shushi-lua.tex
%              または VS Code の「LuaLaTeX (kian-lua)」レシピ
%  プレースホルダー: 【...】を実際の内容に置き換えてください
%  異体字: \ivs{文字} で IPAmj 明朝字形に切り替え（例: \ivs{角󠄁}）
%  ※ \SecA{} 等の {} 内に % を書くとコメントになりエラーが出ます
%  見出し番号: \SecZ=番号なし \SecA=第１ \SecB=１ \SecC=⑴ \SecD=ア \SecE=(ア)
% ============================================================

\documentclass{ltjsarticle}
\usepackage{kian-lua}
% \usepackage{kian-ivs} %IVS 異体字（\ivs{文字}）
% \usepackage{kian-otf} %囲み文字拡張
% \usepackage{luatexja-ruby} %ルビ機能

%  【注意】表紙系コマンド（\Tit \Dat \Sen \Idcs \Recs）・\maketitle は
%  置かない。当事者・裁判所はシステムのフォーム入力で処理される。

\begin{document}

% ---------------------------------------------------------- %
%  請求の趣旨
%  \SecZ は番号なし見出し（「第１」等の番号が付かない）
% ---------------------------------------------------------- %

\SecZ{請求の趣旨}

\begin{enumerate}
  \item 被告は原告に対し，金【請求金額】円及びこれに対する
        令和【年】年【月】月【日】日から支払済みまで
        年【利率（例：３）】分の割合による金員を支払え。
  \item 訴訟費用は被告の負担とする。
\end{enumerate}

との判決並びに仮執行宣言を求める。

\vspace{2\zw}

\begin{flushright}
以　上
\end{flushright}

\end{document}
